\documentclass[a4paper]{article}
\usepackage{graphicx}
\usepackage{amsmath,amssymb}
\usepackage{hyperref}
\usepackage{minted}

\title{}
\date{}
\begin{document}
    \maketitle
    Canary testing is a staged rollout strategy in which a new model is deployed to a small percentage of real users, such as 5\% of total traffic. These users receive actual responses generated by the new model.
    
    System performance and safety metrics are closely monitored. If the results meet predefined thresholds, the rollout percentage is gradually increased, for example from 5\% to 20\%, then 50\%, and eventually 100\%.
    
    The key characteristics of canary testing are:
    
    \begin{itemize}
        \item A small subset of users sees the new model’s outputs.
        \item There is limited but controlled user-facing risk.
        \item It enables validation under real usage conditions.
    \end{itemize}
    
    If problems occur, the deployment can be rolled back quickly.

\end{document}